% Figure: Listing 2, constraint hypotheses as expression edges.
%
% Source for fig_code_constraints.pdf / .png. Build with the sibling driver:
%
%     uv run --group manuscript python manuscript/examples/figures/fig_code_constraints.py
%
% The caption is set INSIDE the float and there is no caption paragraph beneath
% the image, matching Algorithm 1 and Listing 1.
%
% These are the bind() calls made by
% manuscript/examples/fecl4/fecl4_constraint_scenarios.py, whose printed table
% is Table 3. If that script changes, this listing is stale until rebuilt.
\documentclass[border=6pt,varwidth=15.2cm]{standalone}

\usepackage[T1]{fontenc}
\usepackage{lmodern}
\usepackage{xcolor}
\usepackage{listings}
\usepackage{caption}

\definecolor{kw}{RGB}{0,90,160}
\definecolor{str}{RGB}{160,60,20}
\definecolor{cmt}{RGB}{110,110,110}

\lstset{
  language=Python,
  basicstyle=\ttfamily\small,
  keywordstyle=\color{kw}\bfseries,
  stringstyle=\color{str},
  commentstyle=\color{cmt}\itshape,
  showstringspaces=false,
  columns=fullflexible,
  keepspaces=true,
  breaklines=false,
  aboveskip=0.6em,
  belowskip=0pt,
  mathescape=true,
}

\captionsetup{labelfont=bf,skip=0.4em,justification=justified,
              singlelinecheck=false,font=small}
\renewcommand{\lstlistingname}{Listing}

\begin{document}
% Each standalone document restarts the listings counter, so the number is set
% explicitly to match the manuscript's registry. Listing 1 is fig_code_compose.
\setcounter{lstlisting}{1}
\begin{minipage}{15.2cm}
\hrule height 0.9pt
\vspace{0.5em}
\captionof{lstlisting}{Three constraint hypotheses, as expression edges on the
same components. Each \texttt{bind} states a relationship the analyst believes
and the fit must then honour; none of them is expressible to
\texttt{curve\_fit}. The last ties a parameter of one spectrum to a parameter of
another, which is what makes scenario D a single joint fit rather than two fits
compared afterwards.}
\vspace{0.2em}
\hrule height 0.4pt
\begin{lstlisting}
# B: the 2p$_{3/2}$:2p$_{1/2}$ edge-jump degeneracy, 2:1
graph = compose(nodes).bind("step_l3.amplitude / 2", to="step_l2.amplitude")

# C: the L2 white line sits one spin-orbit splitting above the L3 white line,
#    the splitting being the fitted separation of the two continuum steps
graph = graph.bind("p1.center + step_l2.center - step_l3.center", to="p7.center")

# D: ...and that splitting is the SAME atom in both complexes, so the d6 white
#    line is displaced by the splitting the d5 spectrum determines
graph = graph.bind("b_p1.center + a_step_l2.center - a_step_l3.center",
                   to="b_p8.center")
\end{lstlisting}
\vspace{0.3em}
\hrule height 0.9pt
\end{minipage}
\end{document}
